%! Absolute Value Functions \& Transformations — Unit 2, Lesson 2.3 — Homework (Answer Key)
\documentclass[10pt]{article}
\usepackage{algebra2-article}
\usepackage{algebra2-key}   % pulls in -boxes; do NOT also load -boxes
\newcommand{\TallMath}[1]{$\displaystyle #1\rule[-1.4em]{0pt}{3.2em}$}
% Coordinate grid: {xmin}{xmax}{ymin}{ymax}
\newcommand{\lgrid}[4]{%
  \draw[step=1, linegray!40, very thin] (#1,#3) grid (#2,#4);
  \draw[->, semithick, charcoal] (#1,0)--(#2,0) node[below right, font=\scriptsize]{$x$};
  \draw[->, semithick, charcoal] (0,#3)--(0,#4) node[left, font=\scriptsize]{$y$};}

\begin{document}

\pageheader{Unit 2, Lesson 2.3}{Homework --- Answer Key}
\namedateperiod

\vspace{0.10in}

\begin{notesbox}{Practice}
Show your work. For every function, remember: vertex form $a|x-h|+k$ has vertex $(h,k)$ and axis
$x=h$; the sign of $a$ tells you up/down, and $|a|$ tells you narrow/wide.

\begin{enumerate}[leftmargin=1.9em, itemsep=11pt]
  \item \textbf{The parent $y=|x|$.} Fill in each feature.
        \\[2pt] Vertex \ans{$(0,0)$}; \quad axis $x={\color{keyred}\mathbf{0}}$; \quad minimum value
        \ans{$0$}; \quad domain \ans{all real numbers}; \quad range \ans{$y\ge 0$}.

  \item \textbf{Match} each equation to its graph (write the letter).
        \begin{center}
        \begin{tikzpicture}[scale=0.4]
          \begin{scope}
            \lgrid{-1}{5}{-1}{4}
            \draw[very thick, forest] (-1,3)--(2,0)--(5,3);
            \node[charcoal, font=\scriptsize] at (2,-2.4){\textbf{W}};
          \end{scope}
          \begin{scope}[xshift=9cm]
            \lgrid{-3}{3}{-3}{3}
            \draw[very thick, forest] (-3,1)--(0,-2)--(3,1);
            \node[charcoal, font=\scriptsize] at (0,-4.6){\textbf{X}};
          \end{scope}
        \end{tikzpicture}\\[4pt]
        \begin{tikzpicture}[scale=0.4]
          \begin{scope}
            \lgrid{-3}{3}{-4}{2}
            \draw[very thick, forest] (-3,-3)--(0,0)--(3,-3);
            \node[charcoal, font=\scriptsize] at (0,-5.6){\textbf{Y}};
          \end{scope}
          \begin{scope}[xshift=9cm]
            \lgrid{-4}{2}{-1}{5}
            \draw[very thick, forest] (-4,4)--(-1,1)--(2,4);
            \node[charcoal, font=\scriptsize] at (-1,-2.4){\textbf{Z}};
          \end{scope}
        \end{tikzpicture}
        \end{center}
        $y=|x-2|$ is \ans{W}. \quad $y=|x|-2$ is \ans{X}. \quad
        $y=-|x|$ is \ans{Y}. \quad $y=|x+1|+1$ is \ans{Z}.

  \item \textbf{Name the features} of each function.
        \begin{center}
        \renewcommand{\arraystretch}{1.4}
        \begin{tabularx}{\linewidth}{@{}l X X X X@{}}
        & \textbf{Vertex} & \textbf{Axis} & \textbf{Opens} & \textbf{Max/Min value} \\
        (a) $y=|x-2|+1$      & \ans{$(2,1)$} & \ans{$x=2$} & \ans{up} & \ans{min $1$} \\
        (b) $y=-|x+3|$        & \ans{$(-3,0)$} & \ans{$x=-3$} & \ans{down} & \ans{max $0$} \\
        (c) $y=3|x|-4$        & \ans{$(0,-4)$} & \ans{$x=0$} & \ans{up} & \ans{min $-4$} \\
        (d) $y=-\tfrac12|x-1|+2$ & \ans{$(1,2)$} & \ans{$x=1$} & \ans{down} & \ans{max $2$} \\
        \end{tabularx}
        \end{center}

  \item \textbf{Describe the transformations} from $y=|x|$.
        \begin{enumerate}[label=(\alph*), leftmargin=1.8em, itemsep=4pt]
          \item $y=|x-5|+2$: \ans{right $5$, up $2$}
          \item $y=-3|x+1|$: \ans{left $1$, reflect over the $x$-axis (open down), vertical stretch by $3$ (narrower)}
        \end{enumerate}

  \item \textbf{Complete the table} for $y=|x-1|+2$, then give the vertex.
        \begin{center}
        \renewcommand{\arraystretch}{1.2}
        \begin{tabular}{|c|c|c|c|c|c|}
        \hline
        $x$ & $-1$ & $0$ & $1$ & $2$ & $3$ \\ \hline
        $y$ & ${\color{keyred}\mathbf{4}}$ & ${\color{keyred}\mathbf{3}}$ & ${\color{keyred}\mathbf{2}}$ & ${\color{keyred}\mathbf{3}}$ & ${\color{keyred}\mathbf{4}}$ \\ \hline
        \end{tabular}
        \end{center}
        Vertex: $({\color{keyred}\mathbf{1}},\,{\color{keyred}\mathbf{2}})$.

  \item \textbf{Model it.} A drone flies in a straight line past a tower; its distance from the tower
        (in meters) is $D(t)=15\,|t-6|$, where $t$ is in seconds.
        \begin{enumerate}[label=(\alph*), leftmargin=1.8em, itemsep=4pt]
          \item Find the vertex and say what it means in context.
                \ans{Vertex $(6,0)$: at $t=6$ s the drone is \emph{at} the tower (distance $0$, closest).}
          \item What is the distance at $t=0$ and at $t=10$?
                \ans{$D(0)=15(6)=90$ m; \; $D(10)=15(4)=60$ m.}
        \end{enumerate}

  \item \textbf{Explain.} How does the sign of $a$ in $y=a|x-h|+k$ tell you whether the vertex is a
        maximum or a minimum?
        \ansline{If $a>0$ the V opens up, so the vertex is the \emph{lowest} point --- a minimum.}
        \ansline{If $a<0$ the V opens down, so the vertex is the \emph{highest} point --- a maximum.}
\end{enumerate}
\end{notesbox}

\begin{extensionbox}
\textbf{Build and reason.}
\begin{enumerate}[label=(\alph*), leftmargin=1.9em, itemsep=6pt]
  \item A V opens \emph{up}, has vertex $(-2,1)$, and passes through $(0,5)$. Find $a$ and write the
        equation in vertex form.
  \item For $y=a|x-h|+k$ with $a<0$, what is the range in terms of $k$? Explain from the picture.
\end{enumerate}
\ansline{(a) $a=\dfrac{5-1}{0-(-2)}=\dfrac{4}{2}=2$, so $y=2|x+2|+1$.}
\ansline{(b) Range $y\le k$: opening down, the vertex $(h,k)$ is the highest point, so no output exceeds $k$.}
\end{extensionbox}

\begin{spiralbox}
\textbf{Coming next --- Lesson 2.4: Piecewise-Defined Functions.} Today's V is secretly \emph{two}
lines pasted at the vertex: $y=-x$ on the left and $y=x$ on the right. Next you'll write functions
built from different rules on different pieces of the domain --- and meet the greatest-integer (step)
function as a second example.
\end{spiralbox}

\begin{teachernote}
Watch item~4(b): $-3|x+1|$ bundles \emph{three} moves (left $1$, reflect, stretch $3$) --- students
often report only the shift. Item~3(d) tests $a<0$ with a fraction: the graph is \emph{wider} and opens
down, so $2$ is a \emph{maximum}. In item~6, tie the vertex $(6,0)$ to ``closest'' because a distance
can't be negative. Extension~(a): $a$ is the slope of the right ray --- rise $4$ over run $2$ gives
$a=2$; a common slip is to forget that a point $2$ units right of the vertex rose $4$, not $1$.
\end{teachernote}

\end{document}
